\documentclass{article}
\usepackage[utf8]{inputenc}
\usepackage[T1]{fontenc}
\usepackage[margin=2cm]{geometry}
\usepackage{hyperref}
\sloppy
\begin{document}
\section*{Historical source}
Encyclopaedia Britannica scan, source row 1945. Parse 3.5 extracted the text. Super 3.5 VL corrected it from an enlarged crop of each block. It can still contain recognition errors.
Page headers, footers, and page numbers are removed.

\subsection*{Untitled continuation (E0171)}
\textit{Continues an entry from an earlier page.}

server hath not any mark to direct him: hence astro-
nomers are obliged to enlighten the cross-hairs artifi-
cially, in order to render them visible, and this with-
out letting the luminous body interfere with the ob-
ject which they intend to view; and, for doing this,
two ways are proposed.

1. The object-glass of the telescope may be obliquely enlightened by placing a candle near to it in an oblique situation, so that its smoke or flame may not interfere with the object. But if the object-glass should happen to be pretty deep in the tube, it cannot be sufficiently enlightened by this means; and besides, if the telescope is above six feet long, there will be a considerable difficulty in throwing a sufficient quantity of light upon the cross hairs.

2. By some an opening is made in the side of the tube near the focus of the object-glass, through which the cross hairs are illuminated by means of a candle. But this method also is attended with inconveniences; for the observer is incommoded by the light being so near his eyes, and the hairs themselves become liable to accidents through their exposed situation. An error also attends this method; which is, that, according to the position of the light illuminating the hairs, they will appear in different situations. For example, when the horizontal hair is enlightened above, we perceive a luminous line which may be taken for the hair itself, and which appears at its upper superficies. On the other hand, when the hair is enlightened underneath, the luminous line will appear at the under surface of it; and the error will be the diameter of the hair, which often amounts to more than six seconds. M. de la Hire, however, found a remedy for this inconvenience. He often observed, that in moon-shine nights, when the weather was a little foggy, the cross hairs were distinctly seen; but when the heavens were serene, they could scarcely be perceived: he therefore covered with a piece of gauze, or fine silk-crape, that end of the tube next the object-glass; and this method succeeded so well, that a link placed at a good distance from the telescope to enlighten the crape, that the cross-hairs distinctly appeared, and the sight of the stars was by no means obscured.

In making solar observations, a smoked glass must be used for preserving the eyes; and which may be thus prepared. Take two equal and well polished round pieces of flat glass; upon the surface of one of which, all round its limb, glue a pasteboard ring: then put the other piece of glass into the smoke of a lamp, taking it several times out, and putting it in again, left the heat should break it, until the smoke be so thick, that the lamp can scarcely be seen through it: but the smoke must not be all over of the same thickness, that so a place may be chosen answering to the splendour of the sun. This being done, the glass, thus blackened, must be glued to the pasteboard ring abovementioned, with its black side next to the other glass, that the smoke may not be rubbed off.

There are two kinds of observations relating to the flares: one is, when they are in the meridian; the other, when they are in vertical circles. If the position of the meridian be known, the quadrant with which the observation is made, must be placed in the plane of that circle, and then the meridian altitudes are easily observed by means of the plumb-line. The meridian

altitude of a star may likewise be had by a pendulum clock, if the exact time of the star's passing by the meridian be known. It must be observed, however, that stars have the same altitude a minute before and after their passing the meridian, if they be not in or near the zenith; but if they be, their altitudes must be taken every minute when they are in or near the meridian, and then their greatest or least altitudes will be those in question. The position of a given vertical circle must be found by the following method. 1. The quadrant and telescope remaining in the same situation wherein it was when the altitude of a star, together with the time of its passage by the intersection of the cross hairs in the focus of the object-glass, was taken, we observe the time when the sun, or some fixed star, whose longitude and latitude is known, arrives to the vertical hair in the telescope; and from thence the position of the said vertical circle will be had, and also the observed star's true place. But if the sun, or some other star, does not pass by the mouth of the tube; and if a meridian line be otherwise well drawn upon a floor, or very level ground, in the place of observation; a plumb-line must be suspended from a fixed place, at about 18 or 20 feet from the quadrant; under which a mark must be made on the floor, in a right line with the plumb-line. You must next put a thin piece of brass or pasteboard very near the object-glass, in the middle of which there is a slit vertically placed, and passing through the centre of the circular figure of the object-glass. Now, by means of this slit, the before-mentioned plumb-line may be perceived through the telescope, which before could not be seen, because of its nearness. Then the plumb-line must be removed and suspended, so that it be perceived in a right line with the vertical hair in the focus of the object-glass, and a point marked on the floor directly under it. And if a right-line be drawn through this point, and that marked under the plumb-line before it was removed, the said line will meet the meridian drawn upon the floor; and so we shall have the position of the vertical circle in which the observed star is, with respect to the meridian, the angle whereof may be measured in assuming known lengths upon the two lines from the point of concourse; for if, through the extremities of the known lengths, a line or base be drawn, we shall have a triangle, whose three sides being known, the angle at the vertex may be found, which will be the angle made by the vertical circle and meridian.

Under the article Quadrant, is shewn the proper method of fixing that instrument exactly in the meridian: but where the observer has no convenience of this kind, it will be proper to use a portable quadrant; by means of which the altitude of a star must be observed a little before its passage over the meridian, every minute, if possible, until its greatest or least altitude be had; by which means, though we have not the true position of the meridian, yet we know the meridian altitude of the star. But although this method is very good, yet if a star passes by the meridian near the zenith, we cannot have its meridian altitude by repeated observations every minute, unless by chance; because in every minute of an hour, the altitude augments 15 minutes of a degree; and in these kind of observations, the inconvenient situation of the observer, the variation of the star's azimuth several de-


\section*{Source record}
\begin{description}
\item[Scan ID] \texttt{britannica:309f765b1ecc83e0a7f7762e44b4a927a7555443:train:1945}
\item[Image file] \texttt{page-1945-d877e2c0e7d6.jpg}
\item[Image SHA-256] \texttt{d877e2c0e7d6979331208cad696389c92af546a9707a21816f1962452815285f}
\item[Dataset] NationalLibraryOfScotland/encyclopaedia\_britannica\_illustrated at revision \texttt{309f765b1ecc83e0a7f7762e44b4a927a7555443} (CC0-1.0)
\item[Extraction] model \texttt{nvidia/nemotron-3.5-parse}, raw reply \texttt{extract/raw/page-1945-d877e2c0e7d6.json}
\item[Raw reply SHA-256] \texttt{0fee7b7f82deb129f197f050634fa80a095f0d6289358bb47e3fc5f0f17f5aa7}
\item[Correction] model \texttt{nvidia/Nemotron-3.5-Super-VL}, block crops, raw replies \texttt{correct/raw/crop\_nothink/page-1945-d877e2c0e7d6-b*.json}; super\_vl: 7
\item[Running header] Part III. U P T I C S. 5605
\end{description}
\end{document}
